% Generated from locale/tr/content/many-valued-logic/three-valued-logics/kleene.tex; do not hand-edit.


\olfileid[tr]{mvl}{thr}{skl}

\olsection{Kleene Mantıkları}

Stephen Kleene, doğruluk değerlerini hesaplama yordamlarının sonuçları olarak
gören bir anlayışla gerekçelendirilmiş iki üç değerli mantık ortaya koydu:
Bir yordam $\True$ ya da $\False$ sonucunu verebilir, fakat sonlanmayabilir
de. Bu durumda karşılık gelen doğruluk değeri tanımsızdır ve $\Undef$
doğruluk değeriyle gösterilir.

Bir $!A$ önermesinin değillemesini hesaplamak için önce $!A$'nın değerini
hesaplar, ardından sonucun karşıtını döndürürsünüz. $!A$'nın hesaplanması
sonlanmazsa yordamın bütünü de sonlanmaz; dolayısıyla $\Undef$'ın
değillemesi $\Undef$'dır.

Bir $!A \land !B$ tümel evetlemesini hesaplamak için iki seçenek vardır:
Önce $!A$, ardından $!B$ hesaplanabilir; her ikisinin sonucu $\True$ ise
sonuç $\True$, değilse $\False$ olur. Hesaplamalardan biri durmazsa yordamın
bütünü de durmaz. Dolayısıyla bu durumda bileşenlerden biri tanımsızsa tümel
evetleme de tanımsızdır. Aynı durum tikel evetleme için de geçerlidir.

Ancak $!A$ ile $!B$'yi paralel olarak değerlendirebilirsek daha iyisini
yapabiliriz. İki yordamdan biri durup $\False$ döndürdüğünde durabiliriz;
çünkü yanıt yanlış olmak zorundadır. Bu durumda bir bileşeni yanlış olan
tümel evetleme, öteki bileşeni tanımsız olsa bile yanlıştır. Benzer biçimde
bir tikel evetlemeyi paralel hesaplarken iki bileşenden birinin yordamı doğru
sonucunu verir vermez durabiliriz; çünkü tikel evetleme doğru olmak
zorundadır. Dolayısıyla bu durumda bileşenlerden biri tanımsız olsa bile
bileşik önermenin sonucunu bilebiliriz. Bu yorumda $\Undef$'ı “tanımsız”
yerine “bilinmeyen” diye okuyabiliriz.

Bu iki yorum Kleene'in güçlü ve zayıf mantıklarını doğurur. Koşullu,
$\lnot !A \lor !B$'ye eşdeğer olarak tanımlanır.

\begin{defn}
\emph{Güçlü Kleene mantığı}~$\LogKs$ şu matrisle tanımlanır:
\begin{enumerate}
  \item $\lnot$, $\land$, $\lor$, $\lif$ bağlaçlarını içeren standart
  önerme dili $\Lang L_0$.
  \item Doğruluk değerleri kümesi $V = \{\True, \Undef, \False\}$.
  \item Tek belirlenmiş değer $\True$'dur; yani $V^+ = \{\True\}$.
  \item Doğruluk fonksiyonları aşağıdaki tablolarla verilir:
  \begin{center}
    \begin{tabular}{c|c} 
      $\tf{\lnot}$ & \\ 
      \hline  
      $\True$ & $\False$ \\ 
      $\Undef$ & $\Undef$ \\
      $\False$ & $\True$ 
    \end{tabular}
    \quad
    \begin{tabular}{c|ccc} 
      $\tf{\land}[\LogKs]$ & $\True$ & $\Undef$ & $\False$ \\ 
      \hline 
      $\True$ & $\True$ & $\Undef$ & $\False$ \\ 
      $\Undef$ & $\Undef$ & $\Undef$ & $\False$\\ 
      $\False$ & $\False$ & $\False$ & $\False$ 
    \end{tabular}
    \\[2ex]
    \begin{tabular}{c|ccc} 
      $\tf{\lor}[\LogKs]$ & $\True$ & $\Undef$ & $\False$ \\ 
      \hline 
      $\True$ & $\True$ & $\True$ & $\True$ \\ 
      $\Undef$ & $\True$ & $\Undef$ & $\Undef$ \\
      $\False$ & $\True$ & $\Undef$ & $\False$ 
    \end{tabular}
    \quad
    \begin{tabular}{c|ccc} 
      $\tf{\lif}[\LogKs]$ & $\True$ & $\Undef$ & $\False$ \\ 
      \hline 
      $\True$ & $\True$ & $\Undef$ & $\False$ \\ 
      $\Undef$ & $\True$ & $\Undef$ & $\Undef$  \\ 
      $\False$ & $\True$ & $\True$ & $\True$ 
    \end{tabular}
  \end{center} 
\end{enumerate}
\end{defn}

\begin{defn}
\emph{Zayıf Kleene mantığı}~$\LogKw$ şu matrisle tanımlanır:
\begin{enumerate}
  \item $\lnot$, $\land$, $\lor$, $\lif$ bağlaçlarını içeren standart
  önerme dili $\Lang L_0$.
  \item Doğruluk değerleri kümesi $V = \{\True, \Undef, \False\}$.
  \item Tek belirlenmiş değer $\True$'dur; yani $V^+ = \{\True\}$.
  \item Doğruluk fonksiyonları aşağıdaki tablolarla verilir:
  \begin{center}
    \begin{tabular}{c|c} 
      $\tf{\lnot}$ & \\ 
      \hline  
      $\True$ & $\False$ \\ 
      $\Undef$ & $\Undef$ \\
      $\False$ & $\True$ 
    \end{tabular}
    \quad
    \begin{tabular}{c|ccc} 
      $\tf{\land}[\LogKw]$ & $\True$ & $\Undef$ & $\False$ \\ 
      \hline 
      $\True$ & $\True$ & $\Undef$ & $\False$ \\ 
      $\Undef$ & $\Undef$ & $\Undef$ & $\Undef$\\ 
      $\False$ & $\False$ & $\Undef$ & $\False$ 
    \end{tabular}
    \\[2ex]
    \begin{tabular}{c|ccc} 
      $\tf{\lor}[\LogKw]$ & $\True$ & $\Undef$ & $\False$ \\ 
      \hline 
      $\True$ & $\True$ & $\Undef$ & $\True$ \\ 
      $\Undef$ & $\Undef$ & $\Undef$ & $\Undef$ \\
      $\False$ & $\True$ & $\Undef$ & $\False$ 
    \end{tabular}
    \quad
    \begin{tabular}{c|ccc} 
      $\tf{\lif}[\LogKw]$ & $\True$ & $\Undef$ & $\False$ \\ 
      \hline 
      $\True$ & $\True$ & $\Undef$ & $\False$ \\ 
      $\Undef$ & $\Undef$ & $\Undef$ & $\Undef$  \\ 
      $\False$ & $\True$ & $\Undef$ & $\True$ 
    \end{tabular}
  \end{center} 
\end{enumerate}
\end{defn}

\begin{prop}
  $\LogKs$ ile $\LogKw$'nin hiç totolojisi yoktur.
\end{prop}

\begin{proof}
  Bütün !!{propositional variable}s~$p$ için
  $\pAssign v(p) = \Undef$ ise herhangi bir $!A$ formülünün doğruluk değeri
  $\pValue v(!A) = \Undef$ olur; çünkü her iki mantıkta da
  \[
    \tf{\lnot}(\Undef) = \tf{\lor}(\Undef, \Undef) = \tf{\land}(\Undef,
  \Undef) = \tf{\lif}(\Undef, \Undef) = \Undef
  \]
  geçerlidir. Her iki mantıkta da $\Undef \notin V^+$ olduğundan bu
  !!{valuation} altında $!A$ belirlenmiş bir değer almaz.
\end{proof}

Zayıf ve güçlü Kleene mantığının hiç totolojisi olmasa da önemsiz olmayan
gerektirme bağıntıları vardır.

\begin{prob}
  Aşağıdaki bağıntılardan hangileri (a) güçlü ve (b) zayıf Kleene mantığında
  geçerlidir? Her biri için bir doğruluk tablosu verin.
  \begin{enumerate}
    \item $p, p \lif q \Entails q$
    \item $p \lor q, \lnot p \Entails q$
    \item $p \land q \Entails p$
    \item $p \Entails p \land p$
    \item $p \Entails p \lor q$
  \end{enumerate}
\end{prob}

Dmitry Bochvar, $\Undef$'ı “anlamsız” diye yorumladı ve paradoksal tümcelerin
$\Undef$ değerini aldığını şart koşarak Yalancı paradoksu gibi paradoksları
çözmek için kullanmaya çalıştı. Esasen ek bağlaçlarla genişletilmiş zayıf
Kleene mantığı olan bir mantık ortaya koydu; bu bağlaçlardan ikisi “dışsal
değilleme” ile “tanımsızdır” işlemcisidir:
\begin{center}
  \begin{tabular}{c|c} 
    $\tf{\sim}$ & \\ 
    \hline  
    $\True$ & $\False$ \\ 
    $\Undef$ & $\True$ \\
    $\False$ & $\True$ 
  \end{tabular}
\quad
  \begin{tabular}{c|c} 
    $\tf{+}$ & \\ 
    \hline  
    $\True$ & $\False$ \\ 
    $\Undef$ & $\True$ \\
    $\False$ & $\False$ 
  \end{tabular}
\end{center}

\begin{prob}
  Bochvar mantığında $\sim$'yi $\lnot$ ve~$+$ cinsinden tanımlayabilir
  misiniz? Yani !!a{formula} bulun; bu formül yalnızca
  !!{propositional variable}~$p$'yi içersin, $\sim$'yi içermesin ve her zaman
  $\mathord{\sim}p$ ile aynı doğruluk değerini alsın. Doğru olduğunuzu
  göstermek için bir doğruluk tablosu verin.
\end{prob}


